% Источники: exam/materials/преподаватель/2026/Моделирование_случайных_величин,_векторов_.pdf; exam/materials/моделирование.pdf, раздел 13.
\section{Моделирование случайных величин и векторов с нормальным законом распределения.}

\subsection*{Метод Бокса–Мюллера (Box–Muller transform)}

По двум независимым равномерным числам $r_1, r_2 \sim U(0,1)$ строятся два независимых стандартных нормальных числа:
\[
  z_1 = \sqrt{-2\ln r_1}\cdot\cos(2\pi r_2), \qquad
  z_2 = \sqrt{-2\ln r_1}\cdot\sin(2\pi r_2).
\]

Для $N(\mu, \sigma^2)$:
\[
  \xi = \mu + \sigma z.
\]

\subsection*{Метод суммирования (ЦПТ)}

По центральной предельной теореме сумма $n$ независимых $U(0,1)$:
\[
  \xi \approx \frac{\sum_{i=1}^n r_i - n/2}{\sqrt{n/12}} \sim N(0,1).
\]
При $n=12$ упрощается:
\[
  \xi \approx \sum_{i=1}^{12} r_i - 6.
\]
Метод прост, но даёт приближённый результат (хвосты распределения обрезаны).

\subsection*{Нормальный случайный вектор}

Нормальный вектор $\boldsymbol{\xi} \sim N(\boldsymbol{\mu}, \Sigma)$ генерируется через \textbf{разложение Холецкого}:
\[
  \Sigma = AA^\top, \qquad \boldsymbol{\xi} = \boldsymbol{\mu} + A\mathbf{z},
\]
где $A$ — нижнетреугольная матрица (результат разложения Холецкого), $\mathbf{z} = (z_1, \ldots, z_m)^\top$ — вектор независимых $N(0,1)$, полученных методом Бокса–Мюллера.

\textbf{Проверка:} $\mathrm{Cov}(\boldsymbol{\xi}) = A\,\mathrm{Cov}(\mathbf{z})\,A^\top = A I A^\top = \Sigma$. \checkmark

\clearpage
